\section{Rectangular grids}
\label{sec:grids}

Rectangular grids form a basic class of indexing posets in
multiparameter persistence theory.  For integers $m,\ell\geq 2$,
let $G_{m,\ell}:=[m]\times[\ell]$ and
$\Lambda_{m,\ell}:=\Lambda_{G_{m,\ell}}$.
Since $G_{m,\ell}\cong G_{\ell,m}$ as posets, it is enough to consider
$m\geq\ell$.  

For $m\geq\ell\geq2$, define a function
\begin{equation}
\label{eq:grid-combinatorial-function}
c(m,\ell)
:=
\min\{2\ell,m+\ell-2\}
=
\begin{cases}
2\ell-2,&m=\ell,\\
2\ell-1,&m=\ell+1,\\
2\ell,&m\geq\ell+2.
\end{cases}
\end{equation}

Our result is the following explicit formula for the global dimension of $\Lambda_{m,\ell}$.

\begin{theorem}
\label{thm:grid-exact-formula}
For $m\geq\ell\geq2$,
\begin{equation}
\label{eq:grid-exact-formula}
\gldim\Lambda_{m,\ell}=c(m,\ell),
\qquad
\intresgldim_\kk G_{m,\ell}=c(m,\ell)-2.
\end{equation}
\end{theorem}


\begin{proof}
By \Cref{thm:global-dimension-formula} and
\Cref{cor:interval-resolution-global-dimension-formula}, it is enough
to prove
$\Omega(G_{m,\ell})=c(m,\ell)$.

We first prove the upper bound.  Let
$(S,C)\in\Sat(G_{m,\ell})$ and put $T=S\cup C$.
Since $S$ is a relative downset and $C$ a relative upset in $T$,
\begin{equation}
\label{eq:grid-saturated-boundary-inclusions}
\Max(C\setminus S)
=
\Max(T\setminus S)
\subseteq
\Max(T),
\qquad
\Min(S\setminus C)
=
\Min(T\setminus C)
\subseteq
\Min(T).
\end{equation}
Each of $\Min(T)$ and $\Max(T)$ is an antichain, so each meets every
row in at most one element.  Hence
\eqref{eq:grid-saturated-boundary-inclusions} gives
$\bar\omega(S,C)\leq2\ell$.

If $|T|=1$, then $\bar\omega(S,C)=0$.  Suppose that $|T|>1$, and let
$[u,u+w-1]\times[v,v+h-1]$ be the smallest rectangular subgrid
containing $T$. 
Then $|\Min(T)|+|\Max(T)|\leq w+h-1$.
Together with \eqref{eq:grid-saturated-boundary-inclusions}, this gives $\bar\omega(S,C)\leq m+\ell-1$.
We claim that the value $m+\ell-1$ cannot occur.  
Suppose otherwise. In this case, all the inequalities in
\[
\bar\omega(S,C)
\leq
|\Min(T)|+|\Max(T)|
\leq
w+h-1
\leq
m+\ell-1
\]
are equalities.
The last equality gives $w=m$ and $h=\ell$.
Equality in the middle inequality forces
$T=\Min(T)\sqcup\Max(T)$.
Finally, equality in the first inequality forces both inclusions in
\eqref{eq:grid-saturated-boundary-inclusions} to be equalities.
Together with the description of $T$ above, this gives
$S=\Min(T)$ and $C=\Max(T)$.
In this case, $S$ and $C$ are connected antichains.
Thus they are both singletons.
It follows that
$\bar\omega(S,C)=2$, contrary to $m+\ell-1\geq3$.
Thus
\begin{equation}
\label{eq:grid-saturated-pair-upper-bound}
\bar\omega(S,C)
\leq
\min\{2\ell,m+\ell-2\}
=
c(m,\ell).
\end{equation}
Since $(S,C)$ was arbitrary,
$\Omega(G_{m,\ell})\leq c(m,\ell)$.

We now attain the bound.
If $(m,\ell)=(2,2)$, let
\begin{equation}
\label{eq:grid-two-by-two-saturated-pair}
S=\{(1,1)\},
\qquad
C=\{(1,1),(2,1),(1,2)\},
\qquad 
Z=C.
\end{equation}
The only interval properly containing $Z$ is $G_{2,2}$, and $Z$ is
not a relative upset in $G_{2,2}$.  Hence
$\mathcal W(S,C)=\{Z\}$, so $(S,C)$ is saturated.  Moreover,
$\bar\omega(S,C)=2=c(2,2)$.


Assume $(m,\ell)\neq(2,2)$.  For $d\in\mathbb Z$, put
$D_d=\{(x,y)\in G_{m,\ell}\mid x+y=d\}$, and set
\begin{equation}
\label{eq:grid-sharp-saturated-pair}
S:=D_{\ell+1}\cup D_{\ell+2},
\quad
C:=D_{\ell+2}\cup D_{\ell+3}, 
\qand 
Z:=S\cup C.
\end{equation}
See \Cref{fig:grid-SC-three-cases}. 
Both $S$ and $C$ are intervals.  
Moreover, $S$ is a relative downset of $Z$ whose upper boundary is
$D_{\ell+3}=\Max(Z)$, while $C$ is a relative upset of $Z$ whose
lower boundary is $D_{\ell+1}=\Min(Z)$.  Hence
$(S,Z)\in\Upsilon^+(G_{m,\ell})$ and
$(Z,C)\in\Upsilon^-(G_{m,\ell})$, so
$Z\in\mathcal W(S,C)$.

We show that $(S,C)$ is saturated.  Let
$R\in\mathcal W(S,C)$.  Then $Z=S\cup C\subseteq R$.
If $(x,y)\in R\setminus Z$, then either $x+y\leq\ell$ or
$x+y\geq\ell+4$.  In the first case there exists
$a\in D_{\ell+1}\subseteq S$ with $(x,y)\leq a$, contradicting
that $S$ is a relative downset in $R$.  In the second case there
exists $b\in D_{\ell+3}\subseteq C$ with $b\leq(x,y)$,
contradicting that $C$ is a relative upset in $R$.  Thus $R=Z$.
Hence $\mathcal W(S,C)=\{Z\}$, so $(S,C)$ is saturated.

Since
$\Min(S\setminus C)=D_{\ell+1}$ and
$\Max(C\setminus S)=D_{\ell+3}$, we have
\begin{equation}
\label{eq:grid-Z-extremal-count}
\bar\omega(S,C)
=
\ell+\min\{\ell,m-2\}
=
c(m,\ell).
\end{equation}
Therefore $\Omega(G_{m,\ell})\geq c(m,\ell)$.

Together with the upper bound, this gives
$\Omega(G_{m,\ell})=c(m,\ell)$.  The two equalities in
\eqref{eq:grid-exact-formula} now follow from
\Cref{thm:global-dimension-formula} and
\Cref{cor:interval-resolution-global-dimension-formula}.
\end{proof}

\begin{figure}[t]
\centering

\definecolor{gridblue}{RGB}{53,132,244}
\definecolor{gridpurple}{RGB}{128,45,190}
\definecolor{gridred}{RGB}{226,91,82}

\tikzset{
  grid edge/.style={draw=gray!55, line width=0.35pt},
  ambient vertex/.style={circle, draw=gray!60, fill=gray!20, inner sep=1.4pt},
  blue vertex/.style={circle, draw=gridblue, fill=gridblue, inner sep=2.3pt},
  purple vertex/.style={circle, draw=gridpurple, fill=gridpurple, inner sep=2.3pt},
  red vertex/.style={circle, draw=gridred, fill=gridred, inner sep=2.3pt},
  panel label/.style={font=\small},
  grid label/.style={font=\small}
}

\newcommand{\DrawAmbientGrid}[2]{%
  \foreach \i in {1,...,#1}{
    \foreach \j in {1,...,#2}{
      \ifnum\i<#1
        \draw[grid edge] (\i,\j) -- (\the\numexpr\i+1\relax,\j);
      \fi
      \ifnum\j<#2
        \draw[grid edge] (\i,\j) -- (\i,\the\numexpr\j+1\relax);
      \fi
    }
  }
  \foreach \i in {1,...,#1}{
    \foreach \j in {1,...,#2}{
      \node[ambient vertex] at (\i,\j) {};
    }
  }
}

\begin{tikzpicture}[x=0.60cm,y=0.60cm]

%--------------------------------
% (a) m = \ell
%--------------------------------
\begin{scope}[shift={(0,0)}]
  \node[panel label] at (3,6.75) {(a) $m=\ell$};
  \node[grid label] at (3,5.95) {$G_{5,5}$};

  \DrawAmbientGrid{5}{5}

  \foreach \p in {(1,5),(2,4),(3,3),(4,2),(5,1)}{
    \node[blue vertex] at \p {};
  }
  \foreach \p in {(2,5),(3,4),(4,3),(5,2)}{
    \node[purple vertex] at \p {};
  }
  \foreach \p in {(3,5),(4,4),(5,3)}{
    \node[red vertex] at \p {};
  }
\end{scope}

%--------------------------------
% (b) m = \ell+1
%--------------------------------
\begin{scope}[shift={(7.8,0)}]
  \node[panel label] at (3.5,6.75) {(b) $m=\ell+1$};
  \node[grid label] at (3.5,5.95) {$G_{6,5}$};

  \DrawAmbientGrid{6}{5}

  \foreach \p in {(1,5),(2,4),(3,3),(4,2),(5,1)}{
    \node[blue vertex] at \p {};
  }
  \foreach \p in {(2,5),(3,4),(4,3),(5,2),(6,1)}{
    \node[purple vertex] at \p {};
  }
  \foreach \p in {(3,5),(4,4),(5,3),(6,2)}{
    \node[red vertex] at \p {};
  }
\end{scope}

%--------------------------------
% (c) m \ge \ell+2
%--------------------------------
\begin{scope}[shift={(16.6,0)}]
  \node[panel label] at (4,6.75) {(c) $m\geq\ell+2$};
  \node[grid label] at (4,5.95) {$G_{7,5}$};

  \DrawAmbientGrid{7}{5}

  \foreach \p in {(1,5),(2,4),(3,3),(4,2),(5,1)}{
    \node[blue vertex] at \p {};
  }
  \foreach \p in {(2,5),(3,4),(4,3),(5,2),(6,1)}{
    \node[purple vertex] at \p {};
  }
  \foreach \p in {(3,5),(4,4),(5,3),(6,2),(7,1)}{
    \node[red vertex] at \p {};
  }
\end{scope}

%--------------------------------
% Legend
%--------------------------------
\begin{scope}[shift={(1,-0.5)}]
  \node[blue vertex] at (0.8,0) {};
  \node[anchor=west,font=\footnotesize] at (1.4,0)
    {$D_{\ell+1}=S\setminus C=\Min(Z)$};

  \node[purple vertex] at (9.1,0) {};
  \node[anchor=west,font=\footnotesize] at (9.7,0)
    {$D_{\ell+2}=S\cap C$};

  \node[red vertex] at (15.3,0) {};
  \node[anchor=west,font=\footnotesize] at (15.9,0)
    {$D_{\ell+3}=C\setminus S=\Max(Z)$};
\end{scope}

\end{tikzpicture}

\caption{The saturated pair $(S,C)$ in
\eqref{eq:grid-sharp-saturated-pair}. The three panels show the
representative cases $(m,\ell)=(5,5)$, $(6,5)$, and $(7,5)$. Here
$D_{\ell+1}=S\setminus C=\Min(Z)$, $D_{\ell+2}=S\cap C$, and
$D_{\ell+3}=C\setminus S=\Max(Z)$, where $Z=S\cup C$.}
\label{fig:grid-SC-three-cases}
\end{figure}







\begin{remark}
\label{rem:AENY-grid-conjectures}
\Cref{thm:grid-exact-formula} resolves, over an arbitrary coefficient field, the two conjectures on
rectangular grids in
\cite[Conjectures~4.11 and~4.12]{AENY}.
For $\ell=2$ and $m\geq4$, it gives
\begin{equation}
\label{eq:AENY-height-two-conjecture}
\intresgldim_\kk G_{m,2}=2,
\end{equation}
which proves the former conjecture.
More generally, for every fixed $\ell\geq2$,
\begin{equation}
\label{eq:AENY-grid-stabilization}
\intresgldim_\kk G_{m,\ell}=2\ell-2
\qquad
(m\geq\ell+2).
\end{equation}
Thus the interval-resolution global dimension stabilizes at
$2\ell-2$ once $m\geq\ell+2$, proving the latter conjecture.
\end{remark}


